\chapter{基于监督学习的智能协议与参数联合选择方法}
\label{cha:ai_selector}

\section{研究动机与问题描述}

在前几章中，我们分别针对ALOHA类与CSMA类网络提出了集成连续传输机制的改进协议，并通过系统仿真分析了其在吞吐量、时延与公平性等方面的性能表现。 然而，不同业务类型和网络负载条件下的最优协议选择问题仍然缺乏系统化的解决方案。  
我们设定业务到达过程 $\Lambda(t)$，其统计特性在观测周期内保持不变，用以描述业务的静态特征。根据业务模式不同，可定义以下典型流量模型：
\begin{itemize}
	\item \textbf{均匀接入模型}：每个节点在每个时隙均匀概率地到达新包；
	\item \textbf{突发接入模型}：节点群体在以小概率在短时间间隔内集中产生新包；
	\item \textbf{周期性接入模型}：节点以固定周期 $T_p$ 发送业务数据。
\end{itemize}

为此，本章提出一种\textbf{基于监督学习的智能协议与参数选择方法}，旨在在多种随机接入协议并存的系统中，根据业务流量和网络状态等环境特征，自动选择最优的协议类型及其对应的协议特有参数配置。  
与强化学习方法不同，本章采用\textbf{监督学习（Supervised Learning）}框架，通过离线构建数据集、训练模型、在线推断三步实现高效决策。训练数据通过系统仿真获得，模型的输入为统一特征向量 $\mathbf{x}$，输出为最优协议类型及其对应的协议特有参数。


在这里，我们考虑使用马尔科夫链模型来统一刻画三种模型。


\subsection{协议与业务场景的统一描述}

系统支持多种随机接入协议，如 Slotted Aloha、Connection-based Aloha (CBA)、SAST、R-Aloha、CRDSA、IRSA 等。  
为实现统一特征空间表示，需要对任务的输入输出特征设计。

\subsubsection{输入参数定义和分类}

系统输入参数分为两大类：
\begin{itemize}
	\item \textbf{通用参数（共享参数）}：反映网络状态与业务环境特征的参数，所有协议均具备。例如：
	\begin{itemize}
		\item $N$：终端节点数；
		\item $\lambda_b$：平均业务到达率；
		\item $\eta_b$：业务类型指示变量（均匀、突发、周期性等）；
		\item $\sigma_n^2$：信道噪声功率；
		\item 其他通用特征（如平均信噪比、链路占用率等）。
	\end{itemize}
	
	\item \textbf{协议特有参数（专属参数）}：仅在部分协议中生效的控制变量。例如：
	\begin{itemize}
		\item $q$：接入概率（ALOHA类协议特有）；
		\item $M$：连续传输长度（CBA协议特有）；
		\item $L$：副本数（IRSA、CRDSA协议特有）；
		\item $F$：帧长度（CRDSA类协议特有）。
	\end{itemize}
\end{itemize}

为实现统一特征空间表示，定义输入向量为：
\[
\mathbf{x} = [\mathbf{x}_\text{common}, \mathbf{x}_\text{proto}],
\]
其中 $\mathbf{x}_\text{common}$ 表示通用参数部分，$\mathbf{x}_\text{proto}$ 表示协议特有参数部分。
对于不适用的协议参数，在输入向量中仍保留相应位置，并设置“有效性指示位”：
\[
x^{(\text{mask})}_k =
\begin{cases}
	1, & \text{该参数在当前样本中有效；} \\
	0, & \text{该参数不适用。}
\end{cases}
\]

\subsubsection{输出向量定义}
与传统仅预测性能指标的做法不同，本章的输出向量定义为：
\[
\mathbf{y} = [p, \mathbf{x}^*_\text{proto}],
\]
其中：
\begin{itemize}
	\item $p$ 表示预测得到的最优协议类型（例如：Aloha, CBA, IRSA等）；
	\item $\mathbf{x}^*_\text{proto}$ 表示该协议下的最优协议特有参数（如 $q^*, M^*, L^*, F^*$ 等）。
\end{itemize}
即，模型的任务是根据当前网络环境特征，\textbf{联合输出协议类型与其最优参数配置}。

\subsection{性能指标与优化目标}

每次仿真输出以下性能指标：
\[
\mathbf{z} = [T, D, F_f],
\]
其中 $T$ 表示系统吞吐量，$D$ 表示平均时延，$F_f$ 表示公平性指标（如Jain指数）。  
静态场景下的优化目标可表述为：
\[
\max_{\text{protocol},\, \mathbf{x}} \; \mathcal{U}(\mathbf{y}) = w_T T - w_D D + w_F F_f,
\]
其中 $\mathcal{U}(\cdot)$ 为综合性能效用函数，$w_T, w_D, w_F$ 为调节权重。  
智能协议选择器旨在学习从输入特征 $\mathbf{x}$ 到最优协议决策的映射关系。

\section{监督学习数据集构建方法}
\label{sec:dataset_construction}

为训练智能协议选择模型，需要建立跨协议、跨场景的统一监督学习数据集。  
核心思想是：将统一的网络参数视为输入特征，通过系统仿真得到性能输出，从而形成 $(\mathbf{x}, \mathbf{y})$ 对应的数据集样本。

\subsection{仿真样本生成机制}
为训练监督学习模型，首先需构建包含多协议场景的训练数据集。数据生成过程如下：
\begin{enumerate}
	\item 在参数空间中对 $\mathbf{x}_\text{common}$ 和 $\mathbf{x}_\text{proto}$ 进行多维采样；
	\item 对每一组参数组合运行系统级仿真，计算对应性能指标 $\mathbf{z} = [T, D, F_f]$；
	\item 对于每一组环境条件，寻找最优协议类型 $p^*$ 及其最优特有参数 $\mathbf{x}^*_\text{proto}$：
	\[
	(p^*, \mathbf{x}^*_\text{proto}) = \arg\max_{p,\, \mathbf{x}_\text{proto}} \mathcal{U}(T, D, F_f),
	\]
	其中 $\mathcal{U}(\cdot)$ 为综合性能效用函数。
\end{enumerate}
最终得到的监督学习训练样本为：
\[
\mathcal{D} = \{(\mathbf{x}_i, \mathbf{y}_i)\}_{i=1}^{M_s},
\]
其中 $\mathbf{x}_i$ 为输入特征向量，$\mathbf{y}_i = [p_i, \mathbf{x}^*_{\text{proto},i}]$ 为标签输出。

\subsection{数据结构与样本格式}
每个样本的结构如下表所示：

\begin{table}[h]
	\centering
	\caption{监督学习样本结构示例}
	\begin{tabular}{c|c|c|c|c|c}
		\hline
		输入特征 & $N$ & $\lambda_b$ & $\eta_b$ & $q$ & $M,L,F$ \\ 
		\hline
		输出标签 & 协议类型 $p$ & \multicolumn{4}{c}{$[q^*, M^*, L^*, F^*]$} \\ 
		\hline
	\end{tabular}
\end{table}

这样设计后，模型能够同时学习协议选择与参数优化的联合映射关系。


\subsection{模型结构与训练流程}

本问题可建模为\textbf{联合分类与回归任务}，即同时预测最优协议类型（分类）及其协议特有参数（回归）。为此，本章采用多任务深度神经网络（Multi-Task DNN）作为核心算法。网络结构包括：
\begin{itemize}
	\item \textbf{输入层}：接收统一特征向量 $\mathbf{x}$，包括通用网络参数及协议特有参数占位；
	\item \textbf{共享隐藏层}：多层全连接层，用于提取输入特征的高阶表示；
	\item \textbf{输出分支}：
	\begin{enumerate}
		\item 协议选择分支（分类任务）：使用 Softmax 激活函数输出协议类别概率；
		\item 协议参数分支（回归任务）：使用线性激活输出连续协议参数值。
	\end{enumerate}
\end{itemize}

\textbf{训练策略}如下：
\begin{enumerate}
	\item \textbf{数据预处理}：
	\begin{itemize}
		\item 对通用参数及协议特有参数进行归一化或标准化处理；
		\item 协议特有参数不可用位置置零，并使用有效性掩码标记。
	\end{itemize}
	\item \textbf{损失函数设计}：
	\[
	\mathcal{L} = \mathcal{L}_\text{cls} + \alpha \, \mathcal{L}_\text{reg},
	\]
	其中：
	\begin{itemize}
		\item 分类损失 \(\mathcal{L}_\text{cls}\) 为交叉熵：
		\[
		\mathcal{L}_\text{cls} = - \sum_i y_i^{(p)} \log \hat{y}_i^{(p)};
		\]
		\item 回归损失 \(\mathcal{L}_\text{reg}\) 为均方误差（MSE）：
		\[
		\mathcal{L}_\text{reg} = \sum_j \big(x^{*(j)}_\text{proto} - \hat{x}^{(j)}_\text{proto}\big)^2;
		\]
		\item \(\alpha\) 为平衡系数，调整分类与回归任务的重要性。
	\end{itemize}
	\item \textbf{优化算法}：使用 Adam 优化器进行梯度下降训练，初始学习率可设为 $10^{-3}$，并可采用学习率衰减或早停策略。
	\item \textbf{训练过程}：
	\begin{enumerate}
		\item 将数据集分为训练集、验证集和测试集；
		\item 在训练集上迭代优化总损失 $\mathcal{L}$；
		\item 在验证集上监控分类准确率及回归误差，调节网络结构和 \(\alpha\) 超参数；
		\item 最终在测试集上评估模型的综合性能。
	\end{enumerate}
\end{enumerate}

通过上述策略，模型能够同时学习协议类型选择与对应参数预测的联合映射关系，实现多协议、多业务场景下的自适应接入优化。

\subsection{在线执行流程}

在系统运行阶段，智能协议选择器执行以下步骤：
\begin{enumerate}
	\item 系统实时监测业务流量及网络状态特征，例如平均到达率 $\hat{\lambda}_b$、活跃节点数 $\hat{N}$、信道条件等；
	\item 构造完整输入特征向量 $\hat{\mathbf{x}}$，包含通用参数、协议特有参数占位及有效性掩码；
	\item 将输入 $\hat{\mathbf{x}}$ 输入训练好的多任务神经网络，预测输出：
	\[
	(\hat{p}, \hat{\mathbf{x}}_\text{proto}) = f_\text{DNN}(\hat{\mathbf{x}});
	\]
	\item 根据预测的协议类型 $\hat{p}$ 和参数 $\hat{\mathbf{x}}_\text{proto}$ 配置系统接入策略，实现动态最优接入。
\end{enumerate}


\section{小结}

本章提出了一种基于监督学习的智能协议与参数联合选择框架，实现了多协议、多业务场景下的统一建模与自动优化。  
通过数据集采集、标签生成与模型训练，系统能够在运行时根据业务状态自动选择最优协议及其参数配置，为随机接入机制的自适应优化提供了智能化解决方案。